\section{Ternary endpoint-state dynamic program}

At any endpoint cut, each repeated class is in exactly one of three states: not started, active, or finished.  A singleton class has only the not-yet-processed and finished states.  Let $P$ be the finished classes, $Q$ the active repeated classes, and let $R=K\setminus(P\cup Q)$ be the not-started classes.  Put $S=P\cup Q$.

The left span is $\sum_{j\in S}U_j$.  The right span contains precisely classes that have not finished, namely $K\setminus P$, and is therefore $\sum_{j\in K\setminus P}U_j$.  With $\rho(X)=\dim(\sum_{j\in X}U_j)$, Grassmann's identity gives
\[
 \lambda(P,S)
 =\dim\left(\sum_{j\in S}U_j\cap\sum_{j\in K\setminus P}U_j\right)
 =\rho(S)+\rho(K\setminus P)-\rho(K),
\]
because the union of the two class sets is $K$.

Legal transitions are exactly: (i) process a singleton class from not-started to finished; (ii) start a repeated class, moving it from not-started to active; or (iii) finish an active repeated class.  Each transition advances at least one endpoint and never reverses an event, so the state graph is acyclic.  Every legal endpoint order induces one source-to-sink path, and every source-to-sink path spells one legal endpoint order.

Define $D(x)$ as the minimum, over source-to-$x$ paths, of the maximum state width encountered along the path.  For an edge $x\to y$ the standard bottleneck relaxation is
\[
 D(y)\leftarrow\min\{D(y),\max(D(x),\lambda(y))\}.
\]
Processing the acyclic graph in any topological order therefore yields at the sink the minimum possible maximum cut width over all legal endpoint orders.  Endpoint compression then identifies this value with the exact path-width optimum over the original occurrence arrangements.
